% Draft addition to the group-tracking methods/results discussion.
% Merge tracking_references.bib into the Overleaf bibliography before use.
% Unequal-bubble collision results are pending and intentionally not claimed.
\subsection{Delayed Support for Fragmented Shock Surfaces}
\label{subsec:delayed_fragment_support}

An uncertain event may reflect a change in the representation of a surface:
several small detected patches can be followed by one coherent group.
Selecting only a few predecessor patches can then give an incomplete
description. A bounded three-state experiment was therefore used to examine
conditional fragment support. Delaying an association until later observations
are available follows a general principle of multiple-hypothesis tracking
\cite{Reid1979}. Linking already formed track segments provides a related
strategy for handling incomplete observations and merge or split events
\cite{Jaqaman2008}. The present surface-based procedure is an original
adaptation of these ideas, rather than an implementation of either complete
published algorithm.

For three saved times $t_0<t_1<t_2$, the procedure first requires a shared
one-to-one continuation from a group at $t_1$ to a group at $t_2$.
The median normal displacement $s$ between compatible samples in these
groups provides a short-window motion estimate. The earlier position of
each sample in the group at $t_1$ is predicted as
\begin{equation}
    \widetilde{\boldsymbol{x}}(t_0)
    =
    \boldsymbol{x}(t_1)
    - \frac{t_1-t_0}{t_2-t_1}\,s\,\hat{\boldsymbol{n}}(t_1).
    \label{eq:backward_normal_prediction}
\end{equation}
This assumes approximately constant normal motion over the window and
does not identify material-point trajectories. Earlier fragments are compared
with the predicted group using distance and signed-normal compatibility.
An attachment requires at least $50\%$ fragment coverage and a score margin
of $0.1$ over competing eligible families or the unmatched alternative.
The union of attached fragments must cover at least $50\%$ of the predicted
group samples. The number of fragments per family is not limited to three.

The experiment compared the adjacent-event prototype, a static family test,
and the backward-predicted family test. The static test retained the later
continuation requirement but omitted the displacement in
Equation~\ref{eq:backward_normal_prediction}. Across seven targeted synthetic
cases and three distance gates, the respective numbers of exact stipulated
correspondences were $6/21$, $9/21$, and $13/21$. The predicted family method
retained six false links when geometric support was evaluated as physical
ancestry in deliberately indistinguishable-observation controls. These
failures demonstrate that supported family membership must remain distinct
from an established physical predecessor relationship.

In the saved two-bubble windows, the method associated the ten inner patches
at $64^3$ with two families of five, and the fourteen patches at $128^3$ with
two families of seven. At the narrower distance gate, the unions covered
approximately $89.5\%$ and $91.7\%$ of their respective predicted groups;
coverage reached $100\%$ at the wider gate. Family membership was identical
between the static and predicted methods at both gates. Thus, these windows
support the unrestricted fragment-family description, but do not demonstrate
an additional membership benefit from the normal-motion estimate. The
conditional links were retained separately from the existing event graph.

For spatial inspection, the same sample positions were also rendered as
small disks perpendicular to their saved normals, inspired by surfel
representations \cite{Pfister2000}. Each disk had radius $0.45\,\Delta x$.
This display added no connectivity between samples and did not alter the
physical normals, Mach numbers, or surface-integrated measurements. Gaps
remained visible; the rendering did not constitute a closed-surface
reconstruction. Validation with unequal wind strengths remains a subsequent
experiment.
